\chapter{Planteamiento del problema}

En el ámbito de la matemática computacional y el análisis de datos
masivos, el estudio de conjuntos de puntos en espacios de alta
dimensión, $\mathbb{R}^d$ con $d$ elevado, presenta desafíos teóricos
y prácticos significativos.
Gran parte de la metodología clásica en Análisis Multivariante se
fundamenta en la hipótesis de normalidad y en el uso de métricas
euclídeas para la clasificación y el aprendizaje automático.
Sin embargo, la eficacia de estos métodos se ve comprometida por la
denominada \emph{maldición de la dimensionalidad}~\cite{Bellman1958},
fenómeno que describe cómo el crecimiento exponencial del volumen del
espacio diluye la noción de cercanía, afectando la operatividad de
algoritmos basados en distancias, como los $m$-vecinos más próximos o
el Análisis Discriminante Lineal (LDA).
Como respuesta a estas limitaciones, surge el Teorema de
Johnson-Lindenstrauss como una herramienta fundamental para la
reducción de la dimensionalidad.
Este resultado garantiza que es posible proyectar un conjunto finito
de puntos a un espacio de dimensión mucho menor preservando, casi de
manera isométrica, las distancias relativas entre ellos.

\section{Marco formal y objetivos}

Para el desarrollo de este trabajo, se fija un espacio de
probabilidad $\left(\Omega,\Sigma,\mathbb{P}\right)$ y se considera
un conjunto $\mathcal{X}=\{\symbf{x}_{1},\dotsc,\symbf{x}_{n}\}$
compuesto por $n$ puntos en el espacio euclídeo $\mathbb{R}^{d}$.

Definimos una matriz aleatoria $U\in\mathbb{R}^{k\times d}$ de orden
$k \times d$, donde cada entrada $u_{ij}$ es una variable aleatoria
independiente e idénticamente distribuida que sigue una distribución
normal estándar, $u_{ij}\sim\mathcal{N}\left(0,1\right)$.
El objetivo central de este trabajo es describir y demostrar la
validez del Teorema de
Johnson-Lindenstrauss~\ref{thm:johnson-lindenstrauss}, el cual
establece que para cualquier conjunto de $n$ puntos, existe una
proyección lineal hacia un espacio de dimensión
$k\approx\mathcal{O}\left(\varepsilon^{-2}\log n\right)$ que preserva
las distancias euclídeas dentro de un factor de distorsión
$\left(1\pm\varepsilon\right)$.
% Considerando un espacio de probabilidad, donde $\Omega$ es un
% espacio muestral, $\Sigma$ es una $\sigma$-álgebra y $\mathbb{P}$
% es una medida de probabilidad.

% $U$ es la matriz aleatoria de orden $k\times d$, donde cada elemento
% de dicha matriz sigue una distribución normal
% $\mathcal{N}\left(0,1\right)$.
% $\symbf{x}_{1},\dotsc,\symbf{x}_{n}$ son $n$ puntos que están en el
% espacio $\mathbb{R}^{d}$.
Nuestro objetivo es describir y explicar la utilidad del teorema de
Johnson-Lindenstrauss en el contexto de tratamiento de datos de alta
dimensión.
\begin{figure}[ht!]
	\centering
	\includegraphics[width=.55\paperwidth]{jl.png}
	\caption{Esquema de la proyección aleatoria de
		Johnson-Lindenstrauss: reducción de la dimensionalidad
		preservando la estructura métrica de los datos.}
\end{figure}

\section{Estructura del trabajo}

Partiendo de los fundamentos de concentración de medida discutidos
previamente en el curso, este trabajo se estructura de la siguiente
manera:

\begin{itemize}
	\item

	      El Capítulo 2 aborda el desarrollo matemático central,
	      estableciendo el marco formal de las proyecciones aleatorias
	      y proporcionando la demostración detallada del Teorema de
	      Johnson-Lindenstrauss.

	\item

	      El Capítulo 3 se dedica a los ejercicios del libro de
	      texto~\cite{Wegner2024} y experimentos numéricos, donde se
	      resuelven problemas teóricos complementarios y se realiza una
	      validación experimental del teorema mediante simulaciones en
	      Python.
\end{itemize}

Finalmente, el trabajo se completa con un apéndice que recopila
resultados adicionales sobre distribuciones normales multivariantes y
los programas utilizados para las simulaciones, permitiendo validar
empíricamente la robustez de las proyecciones aleatorias en el
tratamiento de datos de alta dimensión.

% Gran parte de los procedimientos clásicos en Análisis Multivariante
% de Datos (Peña 2002; Izenman 2008) se basan en la hipótesis de
% normalidad multivariante.
% Estos procedimientos basados en normalidad son pioneros en
% Clasificación, tanto supervisada como no supervisada, y forman parte
% de pilares metodológicos sobre los que se asienta la Inteligencia
% Artificial, y en particular el Aprendizaje Automático o
% ``Machine Learning'' (Hastie, Tibshirani y Friedman 2009; Bishop
% 2006; James et al. 2014; Shalev-Shwartz y Ben-David 2014).
% Así, por ejemplo, el Análisis Discriminante Lineal (LDA por Linear
% Discriminant Analysis) supone que las observaciones a clasificar han
% sido generadas por distribuciones normales multivariantes con
% matrices de varianzas-covarianzas iguales.
% Este tipo de suposiciones permite clasificar nuevas observaciones
% mediante la estimación de los parámetros asociados a las
% distribuciones normales multivariantes que modelan estas clases.
% Si bien estos métodos gozan de buenas propiedades de optimalidad
% cuando se verifican esas hipótesis de normalidad, se ha comprobado
% que estos métodos se enfrentan a serias dificultades cuando la
% dimensión $d$, por ejemplo, el número de variables discriminantes, crece.
% En este punto, cabe notar que las situaciones de alta dimensionalidad
% son más bien la regla que la excepción en la Ciencia de Datos
% moderna, debido a la omnipresencia de este tipo de datos (registro de
% señales de todo tipo, monitorización de imágenes y vídeo, datos no
% estructurados traducidos a formatos tabulares de alta
% dimensionalidad, \dots).

% Una de las dificultades que aparece al trabajar con la distribución
% normal multivariante en altas dimensiones es el gran número de
% parámetros asociados a todas las covarianzas entre variables a estimar.
% Esto hace que la distancia de Mahalanobis, como forma de medir
% ``proximidades'', naturalmente asociada a la distribución normal multivariante,
% deje de ser operativa a nivel práctico cuando la dimensión $d$ sea elevada.
% Suponer igualdad de varianzas y que las correlaciones entre variables
% son nulas simplifica el uso de la distancia de Mahalanobis al uso de
% distancias euclídeas.
% Así, medir proximidades usando la simple distancia euclídea subyace
% en muchos e importantes métodos básicos en Ciencia de Datos,
% especialmente en métodos que podemos considerar de carácter
% ``no paramétrico''.
% Un ejemplo bien conocido es el método de los $m$-vecinos, que asigna
% observaciones a clases según la clase mayoritaria entre las $m$
% observaciones más cercanas a la que se desea clasificar, dentro del
% espacio euclídeo.

% Desgraciadamente, incluso trabajando en el caso simplificado del uso
% de distancias euclídeas, gran parte de los métodos habituales en
% Ciencia de Datos resultan afectados por un bien conocido fenómeno de
% ``espacio vacío'' que aparece al incrementarse la dimensionalidad de
% los problemas tratados.
% Este fenómeno de espacio vacío suele conocerse como la
% ``maldición de la dimensionalidad'', término introducido por Richard
% Bellman en Bellman (1958), y que, de forma simplificada, es
% atribuible al crecimiento exponencial del volumen del espacio en el
% que se representan los datos cuando se añaden dimensiones.
% En este trabajo se formulará detalladamente este fenómeno de espacio
% vacío y se verá cómo el mismo impacta muy negativamente en muchas de
% las intuiciones de ``baja dimensión'' que sustentan muchas técnicas
% de análisis de datos básicas.
% Por ejemplo, en trabajos como Aggarwal, Hinneburg y Keim (2001) se
% muestra cómo, para la clasificación mediante $m$-vecinos más
% próximos en espacios de muy alta dimensión, sucede que los datos
% son tan escasos a nivel ``local'' que la noción de cercanía o
% vecino más próximo pierde todo su significado.
% Adicionalmente, las técnicas basadas en suposiciones de
% normalidad, aun suponiendo estructuras simplificadas de
% varianzas-covarianzas, dejan de discriminar correctamente entre
% observaciones debido a un fenómeno de concentración de medida que
% se tratará también a lo largo de este trabajo.

% No obstante, a pesar de estas dificultades, en este trabajo se verá
% que, bajo ciertas condiciones bastante generales, la alta dimensión
% puede, en ocasiones, facilitar sorprendentemente algunos análisis de datos.
% Por ejemplo, en espacios de alta dimensión, los vectores aleatorios
% asociados al muestreo de distribuciones normales tienden a ser
% prácticamente ortogonales entre sí, y sus normas euclídeas se
% concentran notablemente en torno a sus valores esperados,
% situándose con alta probabilidad alrededor de dichos valores.
% Este fenómeno de concentración de medida en alta dimensionalidad,
% debidamente aprovechado, simplifica algunas tareas estadísticas
% como, por ejemplo, la separación de clases y permite el
% establecimiento del Teorema de Johnson-Lindenstrauss.
% Es por ello que en algunos contextos se hace referencia también al
% concepto de ``blessing'' (o bendición) de la dimensionalidad, en
% oposición a la maldición de la dimensionalidad.

% Un ejemplo de la aplicabilidad práctica de estos fenómenos de
% concentración de medida y la ortogonalidad en espacios de alta
% dimensión es el Teorema de Johnson-Lindenstrauss.
% Este teorema, formalizado por primera vez en Johnson y
% Lindenstrauss (1984), establece que cualquier conjunto de $n$
% puntos en un espacio $\mathbb{R}^d$ puede proyectarse a un espacio
% de dimensión mucho menor sin sacrificar demasiada información en
% términos de distancias euclídeas.
% Esta proyección, realizada mediante una matriz aleatoria de tamaño
% $k \times d$, preserva aproximadamente las distancias euclídeas
% entre todos los pares de puntos del conjunto original.
% A lo largo de este trabajo se desarrollarán todos los resultados
% previos necesarios para llegar a probar ese Teorema de
% Johnson-Lindenstrauss.
% Este resultado justifica teóricamente el uso de técnicas de
% reducción de dimensionalidad que no requieren asumir ninguna
% estructura particular sobre los datos.
% En particular, resulta de gran ayuda en métodos como la búsqueda de
% $m$-vecinos más próximos y otros métodos intrínsecamente basados en
% distancias euclídeas.
% En virtud de este teorema, al reducir la dimensionalidad mediante
% estas proyecciones aleatorias, se logra que las distancias entre
% las observaciones se conserven dentro de márgenes pequeños, lo cual
% permite, por ejemplo, buscar el vecino más cercano en un subespacio
% de dimensión menor.
% De esta forma, se mitiga parcialmente el problema de alta
% dimensionalidad en contextos de alta complejidad en los datos,
% donde otras técnicas típicas de reducción de la dimensionalidad,
% como por ejemplo las Componentes Principales, resultarían
% completamente inviables o muy costosas computacionalmente.

% La estructura de este trabajo es la siguiente. En el Capítulo 2 se
% presenta la maldición de la dimensionalidad, comenzando con
% simulaciones sobre la distribución normal multivariante
% (Sección 2.1) y la distribución uniforme en hipercubos (Sección 2.2).
% En el Capítulo 3 se estudia la distribución uniforme en la bola
% unidad en alta dimensión.
% En el Capítulo 4 se analiza el comportamiento de la distribución
% normal multivariante en espacios de gran dimensión, profundizando
% en las ideas de concentración introducidas en el capítulo anterior.
% Finalmente, en la Sección 5.1 del Capítulo 5 se introduce el
% Teorema de Johnson-Lindenstrauss y en la Sección 5.2 se muestran
% algunas de sus aplicaciones, incluyendo su uso en algoritmos como
% los de $m$ vecinos más próximos, $m$-medias y LDA para
% clasificación y estimación paramétrica en distribuciones normales.

% Los resultados que se presentan a lo largo de este trabajo están
% basados principalmente en Wegner (2024) a la vez que se han
% consultado otras fuentes y elaborado argumentos propios.
% Asimismo, se ha implementado código estadístico que permite
% justificar empíricamente varios de estos resultados y facilitar su
% comprensión y alcance.
